% ---------------------------------------------------------------------------
% Author guideline and sample document for EG publication using LaTeX2e input
% D.Fellner, v1.20, Jan 18, 2023

\documentclass{egpubl}
\usepackage{STAG2025}
 
% --- for  Annual CONFERENCE
% \ConferenceSubmission   % uncomment for Conference submission
% \ConferencePaper        % uncomment for (final) Conference Paper
% \STAR                   % uncomment for STAR contribution
% \Tutorial               % uncomment for Tutorial contribution
% \ShortPresentation      % uncomment for (final) Short Conference Presentation
% \Areas                  % uncomment for Areas contribution
% \Education              % uncomment for Education contribution
% \Poster                 % uncomment for Poster contribution
% \DC                     % uncomment for Doctoral Consortium
%
% --- for  CGF Journal
% \JournalSubmission    % uncomment for submission to Computer Graphics Forum
% \JournalPaper         % uncomment for final version of Journal Paper
%
% --- for  CGF Journal: special issue
% \SpecialIssueSubmission    % uncomment for submission to , special issue
% \SpecialIssuePaper         % uncomment for final version of Computer Graphics Forum, special issue
%                          % EuroVis, SGP, Rendering, PG
% --- for  EG Workshop Proceedings
% \WsSubmission      % uncomment for submission to EG Workshop
\WsPaper           % uncomment for final version of EG Workshop contribution
% \WsSubmissionJoint % for joint events, for example ICAT-EGVE
% \WsPaperJoint      % for joint events, for example ICAT-EGVE
% \Expressive        % for SBIM, CAe, NPAR
% \DigitalHeritagePaper
% \PaperL2P          % for events EG only asks for License to Publish

% --- for EuroVis 
% for full papers use \SpecialIssuePaper
% \STAREurovis   % for EuroVis additional material 
% \EuroVisPoster % for EuroVis additional material 
% \EuroVisShort  % for EuroVis additional material
% \MedicalPrize  % uncomment for Medical Prize (Dirk Bartz) contribution, since 2021 part of EuroVis

% Licences: for CGF Journal (EG conf. full papers and STARs, EuroVis conf. full papers and STARs, SR, SGP, PG)
% please choose the correct license
%\CGFStandardLicense
%\CGFccby
%\CGFccbync
%\CGFccbyncnd

% !! *please* don't change anything above
% !! unless you REALLY know what you are doing
% ------------------------------------------------------------------------
\usepackage[T1]{fontenc}
\usepackage{dfadobe}  

%\usepackage{placeins}

%\usepackage{cite}  % comment out for biblatex with backend=biber 
% ---------------------------
%\biberVersion
\BibtexOrBiblatex
%\usepackage[backend=biber,bibstyle=EG,citestyle=alphabetic,backref=true]{biblatex} 
%\addbibresource{egbibsample.bib}
% ---------------------------  
\electronicVersion
\PrintedOrElectronic

% for including postscript figures
% mind: package option 'draft' will replace PS figure by a filename within a frame
\ifpdf \usepackage[pdftex]{graphicx} \pdfcompresslevel=9
\else \usepackage[dvips]{graphicx} \fi

\usepackage{egweblnk} 
% end of prologue

% ---------------------------------------------------------------------
% EG author guidelines plus sample file for EG publication using LaTeX2e input
% D.Fellner, v2.04, Dec 14, 2023

\usepackage[ruled, vlined]{algorithm2e} % For algorithms
\LinesNumbered
\newcommand\mycommfont[1]{{#1}}
\SetCommentSty{mycommfont}
\renewcommand{\algorithmcfname}{Algorithm}


\usepackage{paralist}
\usepackage{wrapfig}
\usepackage{amsfonts}

\usepackage{dblfloatfix}
\usepackage{booktabs}
\newcommand{\colrule}{\midrule}
\usepackage[percent]{overpic}

\usepackage[normalem]{ulem}

\newcommand\N{{\mathbb N}}
\newcommand\Q{{\mathbb Q}}
\newcommand\R{{\mathbb R}}
%\newcommand\C{{\mathbb C}}
\newcommand\Z{{\mathbb Z}}
\newcommand\M{{\mathcal M}}
\newcommand\C{{\mathcal C}}
\newcommand{\be}{\begin{equation}}
\newcommand{\ee}{\end{equation}}

\newcommand{\inprod}[2]{{\langle 
  #1 \hspace{0.08cm}, #2 \rangle
  }}
%\theoremstyle{plain}
\newtheorem{theorem}{Theorem}
\newtheorem{prop}{Proposition}[section]
\newtheorem{definition}{Definition}[section]
\newtheorem{lemma}{Lemma}[section]
\newtheorem{example}{Example}[section]
\newtheorem{rem}{Remark}[section]


\usepackage{xspace}
\usepackage{csquotes}
\usepackage{amsmath}
\usepackage{multirow, bigdelim}
% \usetikzlibrary{shapes,decorations.markings,shapes.geometric,positioning,arrows,fadings}
% \usepackage{colortbl}

%%%%% For sectioning:
\newcommand{\parasection} [1] {\textsf{\textbf{#1}.}}
\newcommand{\bs}[1] {\textbf{\textsf{#1}}}

%%%%% For comments:
\newcommand{\oldversion } [1] {}
\newcommand{\ignorethis } [1] {}
\newcommand{\revised} [1] {\textcolor{blue}{ #1}}
%\newcommand{\revised} [1] {#1}
\newcommand{\note       } [1] {\textcolor{red}{\textbf{Note:} {\slshape #1}}}
\newcommand{\smallnote  } [1] {{\small \{#1\}}}
\newcommand{\bignote    } [1] {\begin{quote} \textbf{Note:\ }
                               \slshape #1 \end{quote}}
\newcommand{\urgent } [1] {\textcolor{red}{\textbf{Urgent:}{#1}}}
\newcommand{\issues } [1] {\textcolor{blue}{\textbf{Issue:}{#1}}}
\newcommand{\missing} [1] {\textcolor{green}{\textbf{Missing:}{#1}}}

\newcommand{\gigi} [1] {\textcolor{orange}{\textbf{gg:} {\slshape #1}}}
\newcommand{\enrico } [1] {\textcolor{red}{\textbf{ep:} {\slshape #1}}}
\newcommand{\ciccio } [1] {\textcolor{magenta}{\textbf{cc:} {\slshape #1}}}

\newcommand{\fixme  } [1] {\textcolor{red}{\textbf{#1}}}
\newcommand{\fillit } [1] {\textcolor{cyan}{\textbf{fil:} {\slshape #1}}}
\newcommand{\maybe  } [1] {\textcolor{gray}{\textbf{maybe:} {#1}}}


%%%%% For referencing things:
\newcommand{\chapnum    } [1] {\ref{#1}}
\newcommand{\appnum     } [1] {\ref{#1}}
\newcommand{\sectnum    } [1] {\ref{#1}}
\newcommand{\tblnum     } [1] {\ref{#1}}
\newcommand{\fignum     } [1] {\ref{#1}}
\newcommand{\algnum     } [1] {\ref{#1}}
%\newcommand{\eqnnum     } [1] {\mbox{(\ref{#1})}}
\newcommand{\eqnnum     } [1] {\ref{#1}}
\newcommand{\chap       } [1] {Chapter~\chapnum{#1}}
\newcommand{\chaps      } [1] {Chapters~\chapnum{#1}}
\newcommand{\app        } [1] {Appendix~\appnum{#1}}
\newcommand{\apps       } [1] {Appendices~\appnum{#1}}
\newcommand{\sect       } [1] {Sec.~\sectnum{#1}}
\newcommand{\sects      } [1] {Sections~\sectnum{#1}}
\newcommand{\tbl        } [1] {Tab.~\tblnum{#1}}
\newcommand{\tbls       } [1] {Tables~\tblnum{#1}}
\newcommand{\fig        } [1] {Fig.~\fignum{#1}}
%\newcommand{\figure     } [1] {Figure~\fignum{#1}}
\newcommand{\figs       } [1] {Figures~\fignum{#1}}
\newcommand{\alg 		} [1] {Alg.~\algnum{#1}}
\newcommand{\eqn        } [1] {Equation~\eqnnum{#1}}
\newcommand{\eqns       } [1] {Equations~\eqnnum{#1}}

%%%%% For conditionally referencing things
\newcommand{\cfignum}[2]{\IfRefUndefinedExpandable{#1}{#2}{\fignum{#1}}}


%%%%% Latin:
%% \newcommand{\etal       }     {\textit{et~al.}} old; not like ACM style
\newcommand{\etal       }     {{et~al.}}
\newcommand{\apriori    }     {\textit{a~priori}}
\newcommand{\aposteriori}     {\textit{a~posteriori}}
\newcommand{\perse      }     {\textit{per~se}}
\newcommand{\cf         }     {\textit{cf.}}
\newcommand{\eg         }     {{e.g.,}}
\newcommand{\Eg         }     {{E.g.,}}
\newcommand{\ie         }     {{i.e.,}}
\newcommand{\Ie         }     {{I.e.,}}

%%%%% Math symbols:
\newcommand{\Identity   }     {\mat{I}}
\newcommand{\Zero       }     {\mathbf{0}}
\newcommand{\Reals      }     {{\textrm{I\kern-0.18em R}}}
\newcommand{\isdefined  }     {\mbox{\hspace{0.5ex}:=\hspace{0.5ex}}}
%\newcommand{\implies    }     {\Longrightarrow}
\newcommand{\texthalf   }     {\ensuremath{\textstyle\frac{1}{2}}}
\newcommand{\half       }     {\ensuremath{\frac{1}{2}}}
\newcommand{\third      }     {\ensuremath{\frac{1}{3}}}
\newcommand{\fourth      }    {\ensuremath{\frac{1}{4}}}

%%%%% Math modifiers:
%\renewcommand{\vec      } [1] {{\text{\boldmath $\mathbit{#1}$}}}
\renewcommand{\vec      } [1] {\mathbf{#1}}
\newcommand{\mat        } [1] {{\text{\boldmath $\mathbit{#1}$}}}
\newcommand{\Approx     } [1] {\widetilde{#1}}
\newcommand{\change     } [1] {\mbox{{\footnotesize $\Delta$} \kern-3pt}#1}

%%%%% Math functions:
\newcommand{\Order      } [1] {O(#1)}
\newcommand{\set        } [1] {{\lbrace #1 \rbrace}}
\newcommand{\floor      } [1] {{\lfloor #1 \rfloor}}
\newcommand{\ceil       } [1] {{\lceil  #1 \rceil }}
\newcommand{\inverse    } [1] {{#1}^{-1}}
\newcommand{\transpose  } [1] {{#1}^\mathrm{T}}
\newcommand{\invtransp  } [1] {{#1}^{-\mathrm{T}}}


%%%%% Math functions with small (fixed) and large (expandable) forms:
\newcommand{\abs        } [1] {{| #1 |}}
\newcommand{\Abs        } [1] {{\left| #1 \right|}}
\newcommand{\norm       } [1] {{\| #1 \|}}
\newcommand{\Norm       } [1] {{\left\| #1 \right\|}}
\newcommand{\pnorm      } [2] {\norm{#1}_{#2}}
\newcommand{\Pnorm      } [2] {\Norm{#1}_{#2}}
\newcommand{\inner      } [2] {{\langle {#1} \, | \, {#2} \rangle}}
\newcommand{\Inner      } [2] {{\left\langle \begin{array}{@{}c|c@{}}
                               \displaystyle {#1} & \displaystyle {#2}
                               \end{array} \right\rangle}}
%\newcommand{\argmin     } [1] {{\underset{#1}{\operatorname{argmin}}}}


\newcommand{\manif		} {\mathcal M}
\newcommand{\mesh		} [1] {\ensuremath{\textit{M}^{#1}}}
\newcommand{\meshm		} [1] {\ensuremath{\textit{M\hspace{0.01in}}'^{#1}}}
%\newcommand{\ops		} [1] {\ensuremath{\textit{ops}_{#1}}}
%\newcommand{\graph		} [1] {\ensuremath{\textit{graph}_{#1}}}
\newcommand{\meshtopo   }     {\textcolor{red}{topology}}
\newcommand{\med        }     {\emph{mesh edit distance}}
%\newcommand{\VERT       }     {\textsc{vertex}}
%\newcommand{\FACE       }     {\textsc{face}}

\newcommand{\opc        } [2] {{\ensuremath{#1 \leftrightarrow #2}}}
\newcommand{\opctight   } [2] {{\ensuremath{#1\hspace{-2pt}\rightarrow{}\hspace{-2pt}#2}}}
%\newcommand{\opa        } [1] {{\ensuremath{\epsilon \rightarrow #1}}}
%\newcommand{\opd        } [1] {{\ensuremath{#1 \rightarrow \epsilon}}}
%\newcommand{\opn        }     {\ensuremath{\epsilon \rightarrow \epsilon}}

\newcommand{\extalg     }     {Iterative-Greedy}


%\newcommand{\figlbl     } [1] {{\textbf{\textsf{#1}}}}
\newcommand{\figlbl     } [1] {{{\textsf{#1}}}}


\newcommand{\addinsetbox}[3]{\makebox[0pt]{\hspace{-#1}\raisebox{#2}{\fbox{\includegraphics[width=#3]{figures/blank.png}}}}}
\newcommand{\vhtext}[2]{\begin{sideways}\parbox{#1}{\centering#2}\end{sideways}}
\newcommand{\vhtextr}[2]{\begin{turn}{270}\parbox{#1}{\centering#2}\end{turn}}

% a b c d e
% f g h i j
% k l m n o
\newcommand{\includemesh}[2]{\includegraphics[width=#1]{figures/#2}}

\newcommand{\includemesha}[2]{\includegraphics[width=#1,trim=   0px  720px 1536px   0px,clip=true]{figures/#2}}
\newcommand{\includemeshb}[2]{\includegraphics[width=#1,trim= 384px  720px 1152px   0px,clip=true]{figures/#2}}
\newcommand{\includemeshc}[2]{\includegraphics[width=#1,trim= 768px  720px  768px   0px,clip=true]{figures/#2}}
\newcommand{\includemeshd}[2]{\includegraphics[width=#1,trim=1152px  720px  384px   0px,clip=true]{figures/#2}}
\newcommand{\includemeshe}[2]{\includegraphics[width=#1,trim=1536px  720px    0px   0px,clip=true]{figures/#2}}
\newcommand{\includemeshf}[2]{\includegraphics[width=#1,trim=   0px  360px 1536px 360px,clip=true]{figures/#2}}
\newcommand{\includemeshg}[2]{\includegraphics[width=#1,trim= 384px  360px 1152px 360px,clip=true]{figures/#2}}
\newcommand{\includemeshh}[2]{\includegraphics[width=#1,trim= 768px  360px  768px 360px,clip=true]{figures/#2}}
\newcommand{\includemeshi}[2]{\includegraphics[width=#1,trim=1152px  360px  384px 360px,clip=true]{figures/#2}}
\newcommand{\includemeshj}[2]{\includegraphics[width=#1,trim=1536px  360px    0px 360px,clip=true]{figures/#2}}
\newcommand{\includemeshk}[2]{\includegraphics[width=#1,trim=   0px    0px 1536px 720px,clip=true]{figures/#2}}
\newcommand{\includemeshl}[2]{\includegraphics[width=#1,trim= 384px    0px 1152px 720px,clip=true]{figures/#2}}
\newcommand{\includemeshm}[2]{\includegraphics[width=#1,trim= 768px    0px  768px 720px,clip=true]{figures/#2}}
\newcommand{\includemeshn}[2]{\includegraphics[width=#1,trim=1152px    0px  384px 720px,clip=true]{figures/#2}}
\newcommand{\includemesho}[2]{\includegraphics[width=#1,trim=1536px    0px    0px 720px,clip=true]{figures/#2}}

\newcommand{\includemeshat}[6]{\includegraphics[width=#1,trim=   0px+#2  720px+#3 1536px+#4   0px+#5,clip=true]{figures/#6}}
\newcommand{\includemeshbt}[6]{\includegraphics[width=#1,trim= 384px+#2  720px+#3 1152px+#4   0px+#5,clip=true]{figures/#6}}
\newcommand{\includemeshct}[6]{\includegraphics[width=#1,trim= 768px+#2  720px+#3  768px+#4   0px+#5,clip=true]{figures/#6}}
\newcommand{\includemeshdt}[6]{\includegraphics[width=#1,trim=1152px+#2  720px+#3  384px+#4   0px+#5,clip=true]{figures/#6}}
\newcommand{\includemeshet}[6]{\includegraphics[width=#1,trim=1536px+#2  720px+#3    0px+#4   0px+#5,clip=true]{figures/#6}}
\newcommand{\includemeshft}[6]{\includegraphics[width=#1,trim=   0px+#2  360px+#3 1536px+#4 360px+#5,clip=true]{figures/#6}}
\newcommand{\includemeshgt}[6]{\includegraphics[width=#1,trim= 384px+#2  360px+#3 1152px+#4 360px+#5,clip=true]{figures/#6}}
\newcommand{\includemeshht}[6]{\includegraphics[width=#1,trim= 768px+#2  360px+#3  768px+#4 360px+#5,clip=true]{figures/#6}}
\newcommand{\includemeshit}[6]{\includegraphics[width=#1,trim=1152px+#2  360px+#3  384px+#4 360px+#5,clip=true]{figures/#6}}
\newcommand{\includemeshjt}[6]{\includegraphics[width=#1,trim=1536px+#2  360px+#3    0px+#4 360px+#5,clip=true]{figures/#6}}
\newcommand{\includemeshkt}[6]{\includegraphics[width=#1,trim=   0px+#2    0px+#3 1536px+#4 720px+#5,clip=true]{figures/#6}}
\newcommand{\includemeshlt}[6]{\includegraphics[width=#1,trim= 384px+#2    0px+#3 1152px+#4 720px+#5,clip=true]{figures/#6}}
\newcommand{\includemeshmt}[6]{\includegraphics[width=#1,trim= 768px+#2    0px+#3  768px+#4 720px+#5,clip=true]{figures/#6}}
\newcommand{\includemeshnt}[6]{\includegraphics[width=#1,trim=1152px+#2    0px+#3  384px+#4 720px+#5,clip=true]{figures/#6}}
\newcommand{\includemeshot}[6]{\includegraphics[width=#1,trim=1536px+#2    0px+#3    0px+#4 720px+#5,clip=true]{figures/#6}}

%%%%%% PAPER %%%%%%%%

\newcommand{\AppName}{\emph{SceneGit}\xspace}


\renewcommand{\topfraction}{0.9}    % max fraction of floats at top
\setcounter{topnumber}{2}
\setcounter{totalnumber}{4}     % 2 may work better
\renewcommand{\textfraction}{0.07}  % allow minimal text w. figs
\renewcommand{\floatpagefraction}{0.7}  % require fuller float pages
\renewcommand{\dblfloatpagefraction}{0.7}   % require fuller float pages
\renewcommand{\thefootnote}{\fnsymbol{footnote}}

\DeclareMathOperator*{\argmax}{arg\,max}
\DeclareMathOperator*{\argmin}{arg\,min}
%\DeclareMathOperator*{\avg}{avg}
%\DeclareMathOperator*{\area}{area}
%\DeclareMathOperator*{\mindist}{min-dist}

\newcommand{\quarterpagewidth}{0.8in} %1.0
\newcommand{\thirdpagewidth  }{1.8in} %2.0
\newcommand{\halfpagewidth   }{2.8in} %3.0
\newcommand{\fullpagewidth   }{5.8in} %6.0


\title[Short title]%
      {Persistent Homology vs Scale-space -- A Visual Comparison}
 
\author{Paper XX}
% \author[L. Rocca \& E. Puppo]
%{\parbox{\textwidth}{\centering L.\ Rocca$^{1}$\orcid{0009-0003-3081-2130} 
%and  E.\ Puppo$^{2}$\orcid{0000-0001-9780-5283}
%       }
%      \\
%{\parbox{\textwidth}{\centering $^1$IGAG - CNR, Milan, Italy \hspace{0.5cm}
%$^2$DIBRIS - University of Genoa, Italy }
%}
%}

% end of prologue

\begin{document}

\teaser{
  \setlength\tabcolsep{0pt} % make LaTeX figure out amount of intercol. whitespace
  \noindent
% .png  .png  .png  .png
  \begin{tabular*}{\textwidth}{@{\extracolsep{\fill}}cccc}  
    \includegraphics[width=.24\textwidth]{fig/512x512}
    & \includegraphics[width=.24\textwidth]{fig/512x512} 
    & \includegraphics[width=.24\textwidth]{fig/512x512}  
    & \includegraphics[width=.24\textwidth]{fig/512x512} \\
    A & B & C  & D\\
  \end{tabular*}

  % \includegraphics[width=\linewidth]{eg_new}
  \centering
  \caption{Teaser here}
  \label{fig:teaser}
}

\maketitle
\begin{abstract}
Persistent homology and scale-space analysis are powerful techniques for simplifying the topological structure of scalar data and ranking the importance of a function's critical points. While \emph{persistent homology} works on the \emph{amplitude} of a function by analyzing changes in its sublevel sets, \emph{scale-space} operates on the function's \emph{frequency} through progressive low-pass filtering. Both methods can be used to simplify the \emph{Morse-Smale complex}, which describes the topological structure of a function. Despite their shared objective, it is not well understood under which conditions their classification of topological features significantly differs. We conduct a visual analysis using both synthetic and real-world data to explore the (un)correlation between these two methods, aiming to uncover how the nature of the function's critical points influences the results of each method.

%-------------------------------------------------------------------------
%  ACM CCS 1998
%  (see https://www.acm.org/publications/computing-classification-system/1998)
% \begin{classification} % according to https://www.acm.org/publications/computing-classification-system/1998
% \CCScat{Computer Graphics}{I.3.3}{Picture/Image Generation}{Line and curve generation}
% \end{classification}
%-------------------------------------------------------------------------
%  ACM CCS 2012
  %The tool at \url{http://dl.acm.org/ccs.cfm} can be used to generate
\begin{CCSXML}
<ccs2012>
   <concept>
       <concept_id>10010147.10010371.10010396.10010402</concept_id>
       <concept_desc>Computing methodologies~Shape analysis</concept_desc>
       <concept_significance>500</concept_significance>
       </concept>
   <concept>
       <concept_id>10010147.10010371.10010382.10010383</concept_id>
       <concept_desc>Computing methodologies~Image processing</concept_desc>
       <concept_significance>300</concept_significance>
       </concept>
   % <concept>
   %     <concept_id>10010147.10010371.10010396.10010397</concept_id>
   %     <concept_desc>Computing methodologies~Mesh models</concept_desc>
   %     <concept_significance>300</concept_significance>
   %     </concept>
   <concept>
       <concept_id>10003752.10010061.10010063</concept_id>
       <concept_desc>Theory of computation~Computational geometry</concept_desc>
       <concept_significance>300</concept_significance>
       </concept>
 </ccs2012>
\end{CCSXML}

\ccsdesc[500]{Computing methodologies~Shape analysis}
\ccsdesc[300]{Computing methodologies~Image processing}
% \ccsdesc[300]{Computing methodologies~Mesh models}
\ccsdesc[300]{Theory of computation~Computational geometry}
\printccsdesc
\end{abstract}

%% keywords:
% morse-smale complexes
% morphological data analysis

% !TEX root = STAG25-Morse.tex


\section{Introduction}
Topological data analysis (TDA) finds application in many fields, such as computer graphics \cite{Weinkauf:2009}, scientific visualization \cite{Tierny:2017}, geographical information science \cite{Dey:2017,Rocca:2017fv,Xu:2020}, environmental science \cite{Valsangkar:2019}, genomics \cite{Rabadan_Blumberg_2019}, biomedicine \cite{SKAF2022}, fluid dynamics \cite{Gunther:2012,Sahner:2007}, material science \cite{Gyulassy:2007},  chemistry \cite{Olejniczak:2020}, just to mention a few. 
A recent account of the main techniques in TDA can be found in \cite{Dey_Wang_2022}.

The characterization of functions by their critical points is fundamental to many powerful TDA tools \cite{Biasotti2008, Heine2016, ttk, Dey_Wang_2022}. A prime example is the \emph{Morse-Smale complex} \cite{Smale63}, which captures the morphology of a function $f$ using a geometric graph. The nodes of this graph are the critical points of $f$, and its arcs are the \emph{separatrices} -- the integral lines connecting these critical points.
%, representing ridges and valleys in functions of two variables.
However, a significant challenge arises with real-world, noisy data. Such datasets often contain a multitude of spurious critical points that have little relevance to the function's overall morphology. Consequently, a Morse-Smale complex constructed from the complete set of critical points is frequently so cluttered that it provides minimal or no meaningful insight.

\emph{Persistent homology} \cite{Frosini:1999,Edelsbrunner:2002ve} is a well-established method for ranking the importance of critical points, allowing for the progressive simplification of the Morse-Smale complex \cite{Edelsbrunner:2003dn,Danovaro:2003,Bremer:2004,IURICICH2017,DEY2018,Fugacci:2020,Song:2021}. The method works by analyzing the sublevel sets of a function $f$ to track topological changes that occur at its critical points. This analysis produces a sequence of critical point pairs that are progressively annihilated during topological simplification. Each pair is ranked according to its \emph{persistence}, which corresponds to the difference in function values between the two points.

Broadly speaking, persistent homology operates on the \emph{amplitude} of the signal. 
While this method is very stable to Gaussian noise, it is fragile to impulse noise, which consists of isolated outliers with large amplitudes \cite{reininghaus11}.

An alternative approach to ranking the importance of critical points involves computing the \emph{scale-space} of the input function and analyzing its \emph{deep structure}. 
Originally developed for image processing and low-level computer vision tasks \cite{lindeberg94}, this method applies a diffusion process to the initial function $f$ to create a family of progressively smoother functions. 
The deep structure of this scale-space encodes the morphological information of these functions across different scales. By tracking critical points through these scales, a ranked sequence of pairs is generated; while similar to persistent homology's output, this sequence is distinct.

The analysis performed by this method occurs in the \emph{frequency} domain, making it more robust to impulse noise \cite{Rocca23}. 
However, a challenge arises because critical points drift across the domain as the function undergoes diffusion, which can make reliable tracking difficult \cite{reininghaus11}. 
A continuous, piecewise-linear model of the scale-space can be used to reliably track critical points that vanish together \cite{Rocca:2013}. 
In this case, the ranking of pairs is defined by their resistance to filtering, which is referred to as their \emph{life} within the scale-space.

While both approaches seem effective in simplifying the morphology of a function, it is still not well understood under which conditions their classifications of topological features significantly differ.
In this paper, we present an experimental study that applies both persistent homology and scale-space analysis to a variety of synthetic and real datasets. We use visualization techniques to highlight the specific characteristics and differences in the results obtained from each method.

\enrico{Se viene fuori qualcosa di particolarmente significativo si può anticipare qui.}
%\input{02_related}
% !TEX root = STAG25-Morse.tex


\section{Background notions}
\label{sec:basics}
%\paragraph*{Morse-Smale theory.}
Let ${\mathcal M}$ be a smooth manifold of dimension $d$. 
A smooth function $f:{\mathcal M}\longrightarrow\mathbb{R}$ is said to be a Morse function if all its critical points are isolated; this is equivalent to say that its Riemannian Hessian does not vanish at critical points. 
In the following, we will stick to $d=2$; the theory holds for higher dimensions too, but this is out of the scope of this work. 

%Let $p\in {\mathcal M}$ be a minimum of $f$; we define the \emph{unstable submanifold} (a.k.a, \emph{basin}) of $p$ as the locus of points of ${\mathcal M}$ that lie on integral curves of $f$ emanating from $p$; for $d=2$, each such region is bounded by a set of \emph{separatrices} that are integral curves connecting maxima and saddles. The unstable manifolds form a partition of ${\mathcal M}$.
%Similarly, the \emph{stable submanifold} (a.k.a, \emph{mountain}) of a maximum $q$ is the locus of points that lie on integral curves converging at $q$. 
%The stable submanifolds form another partition of ${\mathcal M}$, and each of them is bounded by separatrices that connect minima to saddles.  
%If the two partitions intersect transversally, then their overlay is called a Morse-Smale complex. 
%Hereafter, we will assume $f$ satisfies this property.
%In the bivariate case, the Morse-Smale complex is described by a planar graph, whose edges are the separatices that connect saddles to maxima and saddles to minima.
%See \cite{matsumoto02} for a more thorough formal treatment of this subject. 

\paragraph*{Persistent homology.}
For any real number $a$, the sublevel set $S_a$ of function $f$ isdefined as the set of all points $p\in\mathcal M$ such that $f(p)\leq a$, i.e., $S_a=f^{-1}((-\infty,a])$. 
As we increase the value of $a$, the sublevel sets grow monotonically, forming a nested sequence of spaces, called a \emph{filtration}.
 Persistent homology analyzes the topological features of $f$ by tracking the changes in the topology of these sublevel sets as $a$ increases.
 
 The topology of the sublevel set $S_a$ changes only when the value of $a$ passes through a critical value of the function $f$. 
 In the case of d=2, these topological changes are linked to the three types of critical points: minima, saddles, and maxima.
%The topological changes are characterized by the Betti numbers of the homology groups. 
%The $k$-th Betti number, $\beta_k$, represents the number of $k$-dimensional "holes" or features in a space. 
%Specifically, $\beta_0$ counts the number of connected components, $\beta_1$ counts the number of one-dimensional loops, and so on. 
%As we sweep through the sublevel sets, a critical point will cause a specific change in the Betti numbers:
%\begin{itemize}
%\item At a minimum, a new connected component is "born" ($\beta_0$ increases by one unit).
%\item At a saddle, a connected component merges with another, or a new loop is formed ($|beta_0$ decreases, or $\beta_1$ increases by one unit).
%\item At a maximum, a loop is filled in, or a new void is created ($\beta_1$ decreases by one unit, or $\beta_2$ increases by one unit). 
%If $\mathcal M$ is compact and connected, a void is created only when the highest value is reached.
%\end{itemize}
The central idea of persistent homology is to pair up these topological events. 
A critical point of index $k$ (which is a saddle, or maximum in our 2D case) is said to pair with a critical point of index $k-1$ (a minimum, or a saddle, respectively). 
The formation of these pairs is best understood by tracking the homology classes of the sublevel sets. 
%A feature's existence is represented by a homology class.
\begin{itemize}
\item A minimum (index 0) corresponds to the birth of a connected component (a 0-dimensional homology class). As the sublevel set grows, this component may persist.
\item A saddle (index 1) can either merge two previously distinct connected components, causing one to "die," or it can create a new loop, causing a new homology class to be "born".
\item A maximum (index 2) always causes the death of a feature, such as a loop being filled in.
\end{itemize}
A persistence pair is therefore a pair of critical points $(p_b,p_d)$ where $p_b$ is a birth point and $p_d$ is a death point. 
For example, a minimum $p_m$ gives birth to a connected component, which persists until a saddle point $p_s$ is reached that merges it with an older component, causing its "death." This forms a persistence pair $(p_m,p_s)$.  
Similarly, a saddle $p_{s'}$  can give birth to a loop which persists until a maximum $p_M$ is reached, causing the loop to be filled and "die". 
This forms the pair $(p_{s'},p_M)$.  

The \emph{persistence} of a pair $(p_b,p_d)$ is defined by the difference between the function values at the death and birth points: $f(p_d)-f(p_b)$.
A large persistence value indicates a robust feature, likely part of the underlying structure, while a small value suggests topological noise.
See \cite{Edelsbrunner:2010} for a more thorough formal treatment of this subject. 

 
\paragraph*{Scale-space.}
The \emph{linear scale-space} $F_f(p,t)$ of $f$ is defined as the solution of the \emph{heat equation} 
\[\frac{\partial}{\partial t} F_f = \lambda \Delta F_f,\]
with initial condition $F_f(p,0)=f(p)$, where $\Delta$ denotes the Laplace-Beltrami operator with respect to the space variable $p$, and $\lambda$ is a constant term tuning the speed of diffusion.
So, the scale-space is defined on a $(d+1)$-dimensional domain ${\mathcal M}t={\mathcal M}\times [0,t_{\max}]$, the first $d$ dimensions referring to space. 
The last dimension is conventionally called \emph{scale} or \emph{time}.
In general, the scale-space $F_f$ is obtained through a diffusion process starting at $f$.
If ${\mathcal M}$ is Euclidean, i.e., ${\mathcal M}\subset\mathbb{R}^d$, the scale-space can be obtained equivalently by convolving $f$ with Gaussian kernels of increasing variance. 

A \emph{slice} of the scale-space for a given time $\bar{t}$ is the restriction $f_{\bar t}(\cdot)=F_f(\cdot,\bar t)$.
Tracking the critical points of slices for varying $\bar t$ provides a set of trajectories through space ${\mathcal M}t$, which constitutes the \emph{deep structure} of the scale-space.
Each critical point $p$ is created at a certain time $t^c_p\geq 0$ and follows one such trajectory, describing its displacement through the scales, until a time $t^a_p\leq t_{\max}$.  

If $t^a_p<t_{\max}$, then the trajectory of $p$ ends together with the trajectory of another critical point $q$, and we say that the pair $(p,q)$ is annihilated at $t^a_{p,q}$. 
If $t^c_p=0$, then $p$ is an \emph{original} critical point of the input function, otherwise it is called a \emph{newborn}. 
Similarly, if $t^a_p=t_{\max}$, then $p$ is a \emph{survivor} at the largest scale.
For $t_{\max}$ large enough, only the absolute minimum and the absolute maximum of $f$ survive. 

A creation or annihilation event always involves two critical points; in 2D, they are always a saddle and either a minimum or a maximum.
Newborn points are usually referred to as topological noise, and either ignored or used to extend the life of original points \cite{Rocca:2013}.
In the following, we will only consider pairs formed of original points.
The \emph{life} of a pair of original points $(p,q)$ that vanish together is the value $t^a_{p,q}$ of the time of their annihilation; the life of survivors is $t_{\max}$.
 
See \cite{Florack:2000wg} for a more thorough formal treatment of this subject. 

\enrico{Forse si potrebbe semplificare evitando di parlare dei newborn.}


%\input{03_classify}
%\input{04_method}
%\input{05_results}
%\input{06_discussion}

\bibliographystyle{eg-alpha-doi} 
\bibliography{references}       

\end{document}

